% =====================================================================
% apendice.tex
% Responsabilidad única: documentar la reproducibilidad (entorno, semilla,
% comandos y comprobaciones automáticas) y la arquitectura por capas del
% paquete. Sin resultados.
% =====================================================================
\section{Reproducibilidad y arquitectura}
\label{sec:apendice}

\subsection{Reproducibilidad}

Todo el proyecto se reconstruye desde la carpeta raíz con \codigo{uv} y
Python 3.12. Cada paso aleatorio parte de la semilla \DatSemilla{} así que
dos ejecuciones dan las mismas cifras. Con \codigo{make notebooks} se ejecutan las siete
etapas en orden y se regenera \codigo{informe/generado/} (las macros, el
índice y las tablas), con \codigo{make colab} los notebooks se copian a una
carpeta vacía y se ejecutan allí sin el paquete para probar que se valen por
sí solos y con \codigo{make pdf} (en \codigo{informe/}) se compila este
documento. Por último \codigo{make verificar} corre el guion
\codigo{verificar.py} de la carpeta \codigo{verificacion/}. Ese guion revisa
la regla de capas, que el generador sea determinista, la ortogonalidad del
L16, las fórmulas de Bayes contra un cálculo a mano, que no haya fugas por
momento, que solo la etapa~07 lea la prueba, que las etapas no tengan cifras
escritas a mano y que no falten macros en el informe.

En Google Colab basta con subir \codigo{notebooks/Proyecto\_Final.ipynb} y
ejecutarlo todo porque su primera sección recrea el paquete. Lo probamos con
una pila parecida a la de Colab (pandas 2.2 y scikit-learn 1.6) y con la local
(pandas 3 y scikit-learn 1.9).

\subsection{Arquitectura por capas}

Cada subcarpeta de \codigo{proyecto/} es una capa y un módulo solo importa de
su capa o de una menor (Tabla~\ref{tab:capas}). Las etapas se comunican solo
por archivos.

\begin{table}[!htbp]
\centering
\small
\begin{tabular}{clp{10.4cm}}
\toprule
Capa & Carpeta & Responsabilidad \\
\midrule
0 & \codigo{base} & Configuración, rutas, artefactos entre etapas y custodia de la prueba \\
1 & \codigo{datos} & Contrato de columnas, festivos, carga, mapeo de la exportación y simulación \\
2 & \codigo{preparacion} & Limpieza determinista, integración de tablas y partición temporal \\
3 & \codigo{variables} & Tablas de modelado con el momento de cada variable (sin fugas) \\
4 & \codigo{modelos} & Tuberías de scikit-learn, modelos, rejillas y agrupamiento \\
5 & \codigo{evaluacion} & Validación por bloques, métricas, umbral por costo, bootstrap y Bayes \\
6 & \codigo{cribado} & Catálogo de análisis, cribado rápido, arreglo L16 y corridas de Taguchi \\
7 & \codigo{reporte} & Figuras, componentes HTML y exportación de cifras y tablas a \LaTeX{} \\
8 & \codigo{etapas} & Una función por sección de notebook: Calcula, guarda y registra cifras \\
\bottomrule
\end{tabular}
\caption{Capas del paquete \codigo{proyecto/}: Los guiones de los notebooks
solo narran y llaman funciones de la capa~8.}
\label{tab:capas}
\end{table}
